\chapter{Literature Review}
\label{chap:literature}

\section{Monetary Transmission Theory}

The starting point is the credit-channel view of monetary policy. In the conventional interest-rate channel, a policy innovation affects short rates, the term structure, borrowing costs, asset prices, and expenditure. Credit-channel models add that policy also changes the external finance premium and borrower balance sheets, so that the effect of a policy shock depends on intermediation frictions and collateral conditions \parencite{BernankeGertler1995}. The financial accelerator extends this logic by showing how changes in borrower net worth amplify shocks through agency costs and credit access \parencite{BernankeGertlerGilchrist1999}.

This framework is central for the thesis because it makes housing credit and bank intermediation first-order transmission objects rather than secondary controls. A policy-news shock can affect the price of credit through lending spreads, the quantity of credit through loan stocks and mortgage flows, and collateral-sensitive sectors through house-purchase lending. The mechanism need not imply that all real-economy variables respond with equal strength or persistence. Heterogeneous responses across lending margins, mortgage credit, broad credit, and compensation-linked variables are a natural feature of the theory.

The post-GFC and post-COVID Euro Area environment made these mechanisms more visible. Policy rates reached effective lower bounds, balance sheets expanded, forward guidance became central, and asset-purchase programs affected term premia, liquidity, duration risk, and portfolio rebalancing. In such an environment, monetary transmission is unlikely to be summarized by a single short-rate movement. It is better understood as a set of channel-specific responses across market prices, bank balance sheets, lending conditions, mortgage credit, and income-linked variables.

Financial intermediaries provide the bridge between policy news and the housing-credit result. \textcite{AdrianShin2010} place intermediary leverage and balance-sheet capacity at the core of monetary economics, while \textcite{BorioZhu2012} emphasize the risk-taking channel. These accounts imply that monetary policy can affect the willingness and capacity of intermediaries to bear risk, price loans, and allocate credit. The thesis therefore treats lending spreads, rates, credit aggregates, and house-purchase lending growth as mechanism variables, not merely as auxiliary outcomes.

\section{High-Frequency Identification and Information Effects}

High-frequency identification addresses the timing problem in monetary economics. At monthly or quarterly frequency, observed policy changes are endogenous to inflation, slack, financial stress, and central-bank forecasts. Measuring market movements inside narrow policy-announcement windows helps isolate unexpected policy news \parencite{Kuttner2001,GertlerKaradi2015}. For the Euro Area, \textcite{Altavilla2019} show that ECB announcements contain several dimensions of news, including target-rate, timing, forward-guidance, and asset-purchase components.

The same literature also explains why interpretation must be cautious. Announcement-window surprises are not automatically pure policy shocks. \textcite{NakamuraSteinsson2018} emphasize that monetary announcements can reveal central-bank information about future fundamentals. \textcite{JarocinskiKaradi2020} develop this insight by distinguishing monetary policy shocks from central-bank information shocks. This matters directly for the thesis because an easing surprise can coincide with adverse macroeconomic information if the central bank signals weaker growth or financial stress.

The DAX result makes this issue unavoidable. A positive easing surprise is followed by a negative real DAX response at medium and longer horizons in the baseline estimates. That sign is inconsistent with a simple asset-appreciation interpretation. The correct reading is that the equity response may combine accommodative repricing with information about future growth, expected earnings, or risk premia. The thesis therefore interprets equity evidence as mixed financial-market evidence under possible information effects, while the housing-credit and lending-condition variables carry the stronger mechanism claim.

External-instrument work reinforces the same discipline. Designs that use policy surprises as instruments require a relevant bridge between the high-frequency surprise and the low-frequency treatment or response object \parencite{MertensRavn2013,StockWatson2018}. Instrument strength diagnostics such as \(F\)-statistics, partial \(R^2\), event concentration, and jackknife sensitivity determine how far the treatment interpretation can be pushed \parencite{StockYogo2005}. The thesis therefore speaks of reduced-form responses to ECB policy-news shocks rather than fully purified structural treatment effects.

\section{Housing-Credit Transmission}

Housing is highly sensitive to monetary policy because it is credit-financed, collateral-intensive, and balance-sheet dependent. Mortgage rates affect user costs and borrowing constraints; collateral values affect credit access; and house-purchase lending connects financial conditions directly to housing-market participation. For these reasons, housing-credit variables deserve separate attention from broad household-credit aggregates.

The housing-credit-cycle literature supports this distinction. \textcite{JordaSchularickTaylor2015} link housing booms, leverage, and financial instability, showing that real estate credit is central to modern financial cycles. \textcite{JordaSchularickTaylor2016} further show that the growth of mortgage finance has transformed the relationship between banking systems, housing, and business cycles. These findings motivate the thesis' focus on house-purchase lending growth as the centerpiece variable.

A key empirical distinction is between stock-based credit growth and flow-based new lending. Stock-based house-purchase credit growth captures cumulative financing conditions, sustained pass-through, and medium-run credit accommodation. Flow-based pure-new mortgage lending captures immediate origination behavior and is more exposed to short-run volatility, sample length, and refinancing or reporting composition. The supervisor's concern that lending growth is strong while pure new loans are weaker is therefore not a contradiction; it is a stock-flow distinction that helps interpret the mechanism.

The thesis' housing-credit claim is consequently narrower than a house-price claim. Quarterly house prices are not interpolated into monthly responses, and the monthly model does not estimate household access outcomes, borrower incidence, or wealth distribution. The identifiable contribution is that house-purchase lending growth responds more clearly and persistently than compensation-linked proxies and broad productive-credit measures.

\section{Financialization and Asset-Price Transmission}

Financialization describes the growing macroeconomic role of financial markets, balance sheets, asset valuations, and market-based intermediation \parencite{Epstein2005,Krippner2005,VanDerZwan2014}. This literature is useful for the thesis, but it must be used in a restrained way. It helps motivate why market prices, liquidity conditions, and intermediary balance sheets may react quickly to policy news; it does not by itself establish welfare, inequality, or redistribution effects.

In a financialized transmission system, policy news can move through discount rates, risk premia, collateral values, portfolio rebalancing, and bank risk-taking. These mechanisms may produce stronger responses in financial and credit variables than in compensation-linked variables, especially when wage adjustment is shaped by bargaining institutions, contracts, expectations, and labor-market frictions. The thesis uses this logic to motivate sectorally differentiated transmission rather than a broad asset-market narrative.

Equity-market responses require special caution. Asset prices are forward-looking and can react immediately, but their signs reflect multiple channels. An easing surprise can raise valuations through discount-rate effects while lowering them through adverse central-bank information about expected growth. The negative DAX response is therefore not hidden or forced into the housing-credit result. It is treated as evidence that financial-market responses are mixed and information-sensitive.

\section{Local Projections and Dynamic Persistence}

Local projections are well suited to the research question because they estimate impulse responses horizon by horizon \parencite{Jorda2005}. A response that appears on impact and disappears by h3 has a different economic meaning from one that remains visible at h6, h12, or h24. Since the thesis is about relative persistence across channels, the path is more informative than a single impact coefficient.

The local-projection framework also clarifies the difference between peak and accumulated evidence. House-purchase lending growth may peak at a medium horizon and remain positive in accumulated terms, while compensation-linked proxies may move at shorter horizons and lose long-horizon precision. This comparison is exactly the reduced-form object the thesis can identify.

LP flexibility must be paired with careful inference. Overlapping horizons create dependence across residuals, and persistent macro-financial variables can make long-horizon inference fragile. HAC covariance estimates and bootstrap checks are therefore part of the evidence rather than an afterthought \parencite{NeweyWest1987}. The thesis' final claim rests on sign, accumulated movement, stability metrics, and uncertainty jointly.

\section{Research Gap}

The literature provides strong tools for studying monetary transmission, but it often emphasizes immediate asset-price movements, broad macro aggregates, or average effects. Less attention is given to whether high-frequency ECB easing shocks transmit more clearly through housing-credit and financial-intermediation variables than through compensation-linked variables.

This thesis fills that gap by combining four elements. First, it uses high-frequency ECB policy-news shocks and monthly local projections. Second, it makes house-purchase lending growth the core empirical contribution. Third, it places lending spreads, credit aggregates, and timing evidence between the policy shock and the housing-credit response. Fourth, it treats market-price results, especially the DAX, as information-sensitive rather than as straightforward evidence of broad asset-price appreciation.

The result is a disciplined macro-financial transmission thesis. It asks where post-COVID-relevant ECB easing shocks remain most visible, shows that the housing-credit channel is the clearest and most persistent response surface, and keeps the interpretation within reduced-form identification boundaries.
